\DTMsavedate{mydate}{2022-06-30}
\maketitle{Digital Electronics}{Electronic Engineering}{\DTMusedate{mydate}}

% ----------------------------------------------------- %
% COURSE INFO
\subsection*{Course Information}\label{course:5111}
\vspace{0ex}%
\noindent\parbox{0.5\textwidth}{%
    \noindent\begin{tabular}{@{}ll}
                 \textsf{Course:}          & 	Digital Electronics \\
                 \textsf{Course Code:}         & 	TEE 2111    \\
                 \textsf{Credits:}    & 	10
    \end{tabular}}
\vspace{2ex}\hrule\vspace{2ex}

% ----------------------------------------------------- %
% INSTRUCTOR INFO
\subsection*{Lecturer Information}
\noindent\parbox{0.5\textwidth}{%
    \noindent\begin{tabular}{@{}ll}
                 \textsf{Name:}         &    Svetlana A. Bebova           \\
% \textsf{Phone:}        &    \phone          \\
\textsf{Email:}        &    \href{mailto:swetlana.bebova@nust.ac.zw}{swetlana.bebova@nust.ac.zw} \\
\textsf{Qualifications:}        &  MSc in Radio Electronic Engineering (Technical University of Sofia)
    \end{tabular}}
\vspace{2ex}\hrule\vspace{2ex}


% ----------------------------------------------------- %
% COURSE SYNOPSIS
\subsection*{Course Synopsis}
This module looks at numerical systems i.e.\ binary, octal, hexadecimal and it also looks at system conversions, mathematical operations in straight and BCD code, logic gates, truth tables, boolean algebra  theorems and K-maps and Minimization of logic expressions.
It also delves into combinational logic applications and design i.e.\ arithmetic circuits, encoders and decoders, code converters, multiplexers and de-multiplexers as well as Flip-Flops
\vspace{2ex} \hrule\vspace{2ex}

% ----------------------------------------------------- %
% LEARNING OBJECTIVES
% \subsection*{Learning Objectives}
% By the end of the course, candidates should be able to:
% \begin{outline}
%    \1      Understand the concept of states as applied to Control Systems.
% \end{outline}
% \vspace{2ex}\hrule\vspace{2ex}

% ----------------------------------------------------- %
% COURSE TOPICS
\subsection*{Topics}
\begin{outline}
    \1 Numerical Systems
    \1 Logic Gates
    \1 Combinational Logic
    \1 Combinational Logic Applications
    \1 Flip-flops
\end{outline}
\vspace{2ex} \hrule\vspace{2ex}

% ----------------------------------------------------- %
% Students Assenssment
\subsection*{Assessment of Students}
\begin{outline}
    \1 Continuous assessment will consist of tests and assignments.
    \1 Final assessment will consist of an examination at the end of the semester.
\end{outline}
\vspace{2ex}\hrule\vspace{2ex}

% ----------------------------------------------------- %
% Grade Evaluation
\subsection*{Course Evaluation and Grading Policies}
Grades are based on:
Continuous assessment (\qty{25}{\percent}),
Final Exam (\qty{75}{\percent})
\vspace{2ex}\hrule\vspace{2ex}
% ----------------------------------------------------- %
% EQUIPMENT
% https://sites.udel.edu/computing-purchases/personal-specs/
% \subsection*{Essential Equipment}

% \vspace{2ex}\hrule\vspace{2ex}
\subsection*{Policies and Other Information}
% Key Text or some Books
\subsection*{Key Text}

 \begin{itemize}
   \item Fundamentals of Digital Electronics, 1998 by B. Paton
   \item Digital Systems Principles and Applications, 10th Edition by R. J. Tocci, N. S. Widmer, G. L. Moss
   \item Theory and Design of Electronic Circuits by E. Tait
 \end{itemize}
\vspace{2ex}\hrule
\vspace{2ex}\hrule\vspace{2ex}